%
% besselj.tex -- Die Bessel-Funktionen J_0, J_1 und J_2 mit den
%                Nullstellen von J_0, Daten aus besselj.dat und
%                besselj-zeros.dat (erzeugt mit besselj.py)
%
% (c) 2026 Gian Cavegn, OST Ostschweizer Fachhochschule
%
\documentclass[tikz]{standalone}
\usepackage{amsmath}
\usepackage{times}
\usepackage{txfonts}
\usepackage{pgfplots}
\pgfplotsset{compat=1.18}
\usetikzlibrary{arrows,intersections,math,calc}
\definecolor{darkred}{rgb}{0.8,0,0}
\definecolor{darkgreen}{rgb}{0,0.6,0}
\begin{document}
\def\skala{1}
\begin{tikzpicture}[>=latex,scale=\skala]
\begin{axis}[
	width=11.4cm,
	height=6cm,
	xmin=0, xmax=20,
	ymin=-0.6, ymax=1.05,
	xtick={0,2,...,20},
	axis lines=middle,
	axis line style={-},
	clip=false,
	grid=both,
	grid style={very thin,gray!25},
	legend pos=north east,
	legend cell align=left,
	tick label style={font=\small}
]
\addplot[thick,blue]       table[x=x,y=J0] {besselj.dat};
\addlegendentry{$J_0(x)$}
\addplot[thick,red]        table[x=x,y=J1] {besselj.dat};
\addlegendentry{$J_1(x)$}
\addplot[thick,darkgreen]  table[x=x,y=J2] {besselj.dat};
\addlegendentry{$J_2(x)$}
\addplot[only marks,mark=*,mark size=1.3pt,blue]
	table[x=x,y=y] {besselj-zeros.dat};
% Achsen über das Gitter hinaus verlängern, wie in den übrigen Abbildungen des Buches
\draw[->] (axis cs:20,0) -- (axis cs:21.2,0) node[below] {$x$};
\draw[->] (axis cs:0,1.05) -- (axis cs:0,1.2);
\end{axis}
\end{tikzpicture}
\end{document}
